\subsection{Adquisición de datos}

Los datos no provienen de campo clínico ni de encuestas: son \textbf{artefactos digitales}
(contratos OpenAPI en YAML/JSON y material de referencia BIAN) que el procedimiento trata como
representaciones documentales comparables. Cada instancia debe identificarse de forma unívoca (p. ej.
\texttt{\{fuente\}\_\{service-domain\}\_\{version\}}).

\begin{table}[H]
  \centering
  \caption{Adquisición de datos --- diseño del experimento (Clase~04 / PI-11)}
  \label{tab:adquisicion-datos}
  \small
  \begin{tabular}{@{}p{0.26\textwidth}p{0.66\textwidth}@{}}
    \toprule
    \textbf{Elemento} & \textbf{En el Grupo~9.2} \\
    \midrule
    Sistema / entorno &
    Sector bancario; arquitecturas API + modelos de referencia BIAN \\
    Objeto de estudio &
    Contrato OpenAPI bancario completo (Cap.~1 §1.4) \\
    Unidad de análisis &
    Par comparable endpoint/\allowbreak schema $\leftrightarrow$ \Behavior{}/\BusinessObject{} \\
    Señal / dato &
    Archivos OpenAPI (YAML/JSON) + extracto del \ServiceDomain{} BIAN \\
    Instrumento de medición &
    Pipeline del modelo (normalización + cálculo S, E, C); \textit{parsing} OpenAPI \\
    Parámetros controlados &
    Versión OpenAPI~3.x; dominio SD; idioma de descripciones; herramienta de \textit{parsing} \\
    Parámetros no controlados &
    Calidad redaccional del contrato; madurez BIAN del emisor; completitud del SD publicado \\
    Factores (Cap.~2) &
    Esquema intermedio; estrategia de matching; técnica semántica; pesos $\alpha,\beta,\gamma$ \\
    Ratios (resultados) &
    \StructuralScore{}, \SemanticScore{}, \CoverageScore{}, \AlignmentScore{}, clasificación \\
    Formato almacenamiento &
    YAML/JSON versionado en repositorio; metadatos CSV/JSON por instancia \\
    Tamaño de muestra &
    Por definir en Tesis~2; ilustración actual: Anexo~F (Payment Initiation) \\
    \bottomrule
  \end{tabular}
\end{table}


\subsection{Criterios de inclusión (muestra empírica)}

Para una validación empírica posterior se considerarán contratos OpenAPI~3.x del sector bancario
emparejados con un \ServiceDomain{} BIAN declarado en la documentación del caso, con:

\begin{itemize}
  \item especificación parseable sin errores de sintaxis;
  \item al menos una operación y un esquema en \texttt{components/schemas};
  \item referencia BIAN disponible en formato legible (catálogo, extracto o documentación oficial);
  \item ausencia de datos personales identificables (Ley N°~29733): preferir ejemplos públicos,
        sintéticos o anonimizados (Anexo~F).
\end{itemize}

La \textbf{población teórica} abarca todas las instancias de medición posibles sobre pares
OpenAPI $\leftrightarrow$ BIAN en ese universo; la \textbf{muestra empírica} será el subconjunto de
contratos/pares seleccionados según representatividad de dominios (pagos, cuentas, préstamos, etc.),
sin fijar un piloto único obligatorio.
